\chapter{BCH编码方案与仿真分析}\label{chap:bch}

\section{BCH低速电文编码设计目标}

本章针对长度为200 bit和300 bit的低速二进制电文，研究Bose--Chaudhuri--Hocquenghem（BCH）码的编码组织、软件实现和信道适应性。此类电文的数据量较小，但通常要求整帧正确交付；只要任一信息比特恢复错误，该帧就不能作为有效业务电文使用。因此，方案评价以帧错误率（Frame Error Rate，FER）为主，以比特错误率（Bit Error Rate，BER）、发送长度、实际码率和软件译码时间为辅助指标。

设计工作围绕两条路线展开。第一条路线把原始电文划分为若干11 bit信息组，各组独立采用BCH(15,11)编码，主要特点是结构简单、单组译码开销小。第二条路线把一帧信息装入较长的BCH母码，通过缩短获得与业务长度相适应的码字，利用较大的整块纠错能力提高可靠性。两条路线采用相同的二进制相移键控（Binary Phase Shift Keying，BPSK）调制和硬判决条件，使比较重点落在编码组织与纠错能力本身。

除加性高斯白噪声（Additive White Gaussian Noise，AWGN）基准外，本章还考察固定多径、频偏造成的帧内相位累积、连续遮挡以及连续突发翻转。研究目标不是给出脱离使用条件的单一“最优码”，而是回答三个工程问题：在相同信息长度下，冗余和纠错能力如何影响可靠性；简单查表译码与长码代数译码需要付出多大软件代价；当错误从分散形态转为局部集中形态时，单纯增强BCH码是否仍然有效。

\section{BCH编码组织方案}

\subsection{分组BCH方案}

分组方案以11 bit为一个信息组，逐组采用可纠正单个错误的BCH(15,11)系统码。若最后一组不足11 bit，则在信息尾部加入已知零作为填充位；接收端译码并拼接各组信息后，再删除填充位，得到原始电文。200 bit电文被划分为19组，加入9个填充位，发送长度为285 bit；300 bit电文被划分为28组，加入8个填充位，发送长度为420 bit。

对每个信息组，系统编码可写为
\begin{equation}
  c(x)=x^{n-k}m(x)+r(x),
\end{equation}
其中，$m(x)$为信息多项式，$r(x)$为$x^{n-k}m(x)$除以生成多项式所得的余式。信息位在码字中保持可见，便于检查编码与译码接口。由于每个子码字只能保证纠正1 bit错误，整帧包含的子块越多，任一子块出现多错并导致整帧错误的机会也越大；另一方面，错误被限制在局部子块内，译码器只需处理15 bit短码字，软件实现直观且耗时较低。

\subsection{整块缩短BCH方案}

整块方案首先选择信息容量不小于原始电文长度的BCH母码，再在母码信息区前端补入已知零。完成母码系统编码后，发送端删除这些已知零对应的位置；接收端在相同位置补回零，再按完整母码进行译码。该过程称为缩短。缩短不改变母码的设计纠错能力$t$，但减少了实际发送位数，实际码率应按
\begin{equation}
  R=\frac{K_{\mathrm{info}}}{N_{\mathrm{tx}}}
\end{equation}
计算，其中$K_{\mathrm{info}}$为原始信息比特数，$N_{\mathrm{tx}}$为实际发送码长。

对于200 bit电文，候选母码包括BCH(255,207)、BCH(511,421)和BCH(511,385)，相应保证纠错能力分别为6、10和14 bit。对于300 bit电文，BCH(511,421)和BCH(511,385)均可采用单个缩短码字；BCH(255,207)的信息容量不足以容纳整帧，因而把电文平均分为两个150 bit子块并分别编码。双块方式仍属于长码代数译码，但每个子块独立承担错误，帧错误率由两个子块的联合恢复结果决定。

图\ref{fig:bch-structure-placeholder}用于后续补充两类编码组织的Visio结构图。图中应明确标出信息分组、填充位、缩短零前缀、校验位以及最终发送序列，便于从接口层面区分“末组补齐”和“母码缩短”这两种码长适配方法。

\begin{figure}[htbp]
  \centering
  \fbox{\begin{minipage}[c][0.21\textheight][c]{0.82\textwidth}
    \centering\Large 待插入：Visio绘制的BCH分组编码与整块缩短编码结构图
  \end{minipage}}
  \caption{BCH分组编码与整块缩短编码结构}
  \label{fig:bch-structure-placeholder}
\end{figure}

\subsection{两类方案结构与参数比较}

表\ref{tab:bch-schemes}给出八种编码方案的业务参数。表中不使用程序内部索引，编码长度即进入BPSK调制器的比特数。“填充”发生在分组方案的信息尾部，“缩短”则是在长母码信息区加入并删除已知零，两者不能混为同一操作。

\begin{table}[htbp]
  \centering
  \scriptsize
  \caption{两类BCH编码方案的结构与参数}
  \label{tab:bch-schemes}
  \setlength{\tabcolsep}{2.2pt}
  \begin{tabular}{c p{0.19\textwidth} c c c c c c p{0.12\textwidth}}
    \toprule
    信息长度 & 编码组织 & 母码 & 块数 & 填充/缩短 & 发送长度 & 码率 & 纠错能力 & 译码方法 \\
    \midrule
    200 & 200 bit分组方案 & (15,11) & 19 & 填充9 & 285 & 0.7018 & 每块1 & 综合查表 \\
    200 & 200 bit整块方案 & (255,207) & 1 & 缩短7 & 248 & 0.8065 & 6 & 代数译码 \\
    200 & 200 bit整块方案 & (511,421) & 1 & 缩短221 & 290 & 0.6897 & 10 & 代数译码 \\
    200 & 200 bit整块方案 & (511,385) & 1 & 缩短185 & 326 & 0.6135 & 14 & 代数译码 \\
    300 & 300 bit分组方案 & (15,11) & 28 & 填充8 & 420 & 0.7143 & 每块1 & 综合查表 \\
    300 & 300 bit双块方案 & (255,207) & 2 & 每块缩短57 & 396 & 0.7576 & 每块6 & 代数译码 \\
    300 & 300 bit整块方案 & (511,421) & 1 & 缩短121 & 390 & 0.7692 & 10 & 代数译码 \\
    300 & 300 bit整块方案 & (511,385) & 1 & 缩短85 & 426 & 0.7042 & 14 & 代数译码 \\
    \bottomrule
  \end{tabular}
\end{table}

200 bit BCH(255,207)整块方案的发送长度最短、码率最高，但纠错能力低于两个511 bit母码方案；BCH(511,385)整块方案以较多校验位换取$t=14$。300 bit BCH(511,421)整块方案只发送390 bit，比双块BCH(255,207)方案还少6 bit，同时把纠错范围扩展到整个码字；BCH(511,385)整块方案进一步把$t$提高到14，但发送长度增至426 bit。由此可见，码长、码率和纠错能力共同变化，不能只根据母码长度判断性能。

\section{BCH编译码实现}

\subsection{系统编码与码长适配}

编码器首先检查输入长度及二进制取值。分组方案补齐末组后逐块计算校验位；整块方案构造“已知零前缀+原始信息比特”的母码信息区，完成系统编码后删除零前缀。接收端执行相反的长度适配，只把去除填充位或缩短前缀后的原始信息区交给统计模块。BCH(255,207)在$\mathrm{GF}(2^8)$上运算，两个511 bit母码在$\mathrm{GF}(2^9)$上运算，有限域表示、生成多项式和码字位序在编码器与译码器之间保持一致。

\subsection{分组BCH查表译码}

对BCH(15,11)，接收子码字先计算综合。综合为零时直接输出信息位；综合非零时，以综合值查询预先建立的单错位置表，翻转对应比特并再次校验。由于$t=1$且码长很短，查表法避免了迭代求错误定位多项式，核心操作只是有限域计算、数组访问和一次候选翻转，适合低复杂度软件或逻辑电路实现。

若一个15 bit子块内出现两个及以上错误，译码结果可能有两种表现：译码器找不到满足条件的错误位置并声明失败，或者接收向量落入另一合法码字的判决区域，译码器虽然声明成功，却恢复出错误的信息。这也是为什么整帧统计不能只读取译码器状态。

\subsection{整块BCH代数译码}

整块译码先补回缩短位置，再计算$2t$个综合；随后采用Berlekamp--Massey（BM）算法求取错误定位多项式，并用Chien搜索遍历母码位置寻找根。仅当根数与定位多项式次数一致、根不落在已删除的缩短前缀中、翻转候选错误后综合重新归零时，译码器才声明成功。最后去除校验区与缩短前缀，恢复原始电文。

代数译码的主要开销来自综合计算、BM迭代和对255或511个母码位置进行Chien搜索。随着$t$和母码长度增加，有限域乘加次数明显高于15 bit子块的查表路径；其收益是能够在整个长码字范围内联合处理更多分散错误。图\ref{fig:bch-decode-placeholder}预留为两类译码流程图，后续Visio绘制应呈现长度适配、综合计算、错误定位、后验校验和信息提取的差异。

\begin{figure}[htbp]
  \centering
  \fbox{\begin{minipage}[c][0.21\textheight][c]{0.82\textwidth}
    \centering\Large 待插入：Visio绘制的分组BCH与整块BCH译码流程图
  \end{minipage}}
  \caption{分组BCH与整块BCH译码流程}
  \label{fig:bch-decode-placeholder}
\end{figure}

\subsection{译码结果判定}

BCH码的保证纠错能力满足
\begin{equation}
  t=\left\lfloor\frac{d_{\min}-1}{2}\right\rfloor ,
\end{equation}
但该式只给出不超过$t$个错误时的保证，并不承诺错误数超过$t$后一定能被检测出来。为避免把误纠计为成功，仿真采用两级判定：第一级记录译码器是否声明完成纠错，第二级将恢复后的原始信息比特与发送电文逐位比较。只有两者一致，才记为实际正确恢复；声明完成但信息仍错误的帧记为误纠，不能形成有效输出的帧记为译码失败。FER最终按完整电文是否正确计算，因此分组方案中任一子块造成的信息错误都会使整帧计错。

\section{AWGN条件下的BCH方案选择}

\subsection{仿真条件}

AWGN比较采用统一BPSK调制、实值高斯噪声和硬判决译码。为了在瀑布区获得足够错误样本，同时控制高信噪比点的运行时间，各信噪比点设置最少帧数、目标错误帧数和最大帧数三重停止条件。表\ref{tab:bch-awgn-settings}列出正文需要的仿真条件，随机种子和原始数据列名等复现信息集中列于附录\ref{app:reproducibility}。

\begin{table}[htbp]
  \centering
  \small
  \caption{AWGN条件下的BCH方案比较设置}
  \label{tab:bch-awgn-settings}
  \begin{tabular}{p{0.22\textwidth}p{0.65\textwidth}}
    \toprule
    项目 & 设置 \\
    \midrule
    输入电文 & 200 bit、300 bit独立等概率二进制序列 \\
    候选编码 & 表\ref{tab:bch-schemes}所列8种分组、单块缩短和双块缩短方案 \\
    调制与信道 & BPSK；实值AWGN；硬判决 \\
    横轴范围与步长 & $E_s/N_0=-3.0103\sim14.9897$ dB，步长0.5 dB \\
    停止条件 & 每点至少1000帧，累计200个错误帧后可停止，每点最多50000帧 \\
    结果判定 & 恢复的原始信息比特与发送电文逐位比较 \\
    主要指标 & FER；同时统计BER和软件译码时间 \\
    \bottomrule
  \end{tabular}
\end{table}

对幅度为$\pm1$的实值BPSK，噪声方差定义对应的横轴统一换算为符号能量与噪声功率谱密度之比$E_s/N_0$。高信噪比点若在50000帧内未观察到错误，图中保留其“零观测”事实，但不把它解释为理论错误率为零，也不人为绘制错误平台。

\subsection{200 bit方案}

图\ref{fig:bch-s1-k200-fer}显示，200 bit BCH(511,385)整块方案在所考察的AWGN范围内可靠性最好，其后依次为BCH(511,421)整块方案、BCH(255,207)整块方案和BCH(15,11)分组方案。为给工程选型提供易比较的量化位置，在相邻两个非零统计点之间对$\log_{10}(\mathrm{FER})$作线性插值。FER为$10^{-2}$时，四种方案所需$E_s/N_0$约为2.95、3.52、4.39和5.98 dB。

\reportfigure[0.86\textwidth]{figures/bch/s1_awgn_k200_fer.png}
  {200 bit电文的AWGN帧错误率}{fig:bch-s1-k200-fer}
  {横轴统一为$E_s/N_0$；零观测点不作为非零对数曲线绘制。}

上述插值只用于描述相邻采样点之间的曲线位置，不等同于理论门限，也没有扩大原始样本量。BCH(511,385)整块方案通过326 bit发送长度和更高代数译码开销取得最靠左的曲线；BCH(255,207)整块方案只发送248 bit，但$t=6$限制了低错误率区域的收益；分组方案虽然单块译码简单，却要连续正确恢复19个子块，因此整帧FER较高。实际选择需要同时考虑链路余量、占用时长和处理资源。

\subsection{300 bit方案}

图\ref{fig:bch-s1-k300-fer}给出300 bit电文结果。FER为$10^{-2}$时，BCH(511,385)整块方案、BCH(511,421)整块方案、双块BCH(255,207)方案和BCH(15,11)分组方案的插值位置约为3.46、4.03、4.27和6.18 dB。390 bit BCH(511,421)整块方案比396 bit双块BCH(255,207)方案略靠左，说明在当前分散噪声条件下，把$t=10$作用于整个码字获得的联合纠错收益，抵消了二者很小的码率差异。

\reportfigure[0.86\textwidth]{figures/bch/s1_awgn_k300_fer.png}
  {300 bit电文的AWGN帧错误率}{fig:bch-s1-k300-fer}
  {各曲线使用相同的调制、噪声与整帧正确性判据。}

426 bit BCH(511,385)整块方案相对390 bit方案把FER为$10^{-2}$的位置改善约0.57 dB，但多发送36 bit，并增加代数译码运算。双块方案的两个子码字各自具有$t=6$，如果错误恰好分散到两块，它可能获得局部并行纠错的机会；然而整帧需要两个子块都正确，且当前软件按顺序完成两次代数译码，所以不能仅按“每块$t=6$”与单块$t=10$直接换算优劣。

\subsection{可靠性、码长与软件译码时间权衡}

图\ref{fig:bch-s1-latency}给出同一批AWGN仿真中的每帧平均软件译码时间。按全部信噪比点累计译码时间除以累计帧数，200 bit四类方案从分组BCH(15,11)、BCH(255,207)、BCH(511,421)到BCH(511,385)约为12.73、20.54、43.66和63.67~$\mu$s/帧；300 bit对应约为18.52、28.60、39.00和50.96~$\mu$s/帧。分组方案虽然需要19或28次子块译码，但每次只做短码综合与查表；长码方案则承担BM迭代和完整母码范围的Chien搜索。

\reportfigure[0.92\textwidth]{figures/bch/s1_decode_latency.png}
  {AWGN条件下各BCH方案的软件译码时间}{fig:bch-s1-latency}
  {结果为本地Windows/MinGW软件运行平台的平均测量值。}

这些时间反映当前处理器、编译器和串行软件实现，不代表现场可编程门阵列或专用集成电路的流水线时延，也不应被视为与平台无关的算法常数。综合可靠性和代价，200 bit场景若强调低复杂度，可采用分组BCH(15,11)；若可靠性优先，则BCH(511,385)整块方案更有优势。300 bit场景中，390 bit BCH(511,421)整块方案在发送长度、AWGN可靠性和软件时间之间形成较均衡的基线，426 bit BCH(511,385)整块方案适用于愿意增加冗余以换取更低FER的场合。

\section{不同信道条件下的BCH适应性}

\subsection{多信道仿真设置}

多信道试验沿用表\ref{tab:bch-schemes}中的八种编码方案，但根据失真机理分别设置横轴和接收处理，不能把不同图中的数值拼成统一排名。表\ref{tab:bch-multichannel-settings}概括各场景。图\ref{fig:bch-multichannel-placeholder}预留为Visio链路图，后续应画出编码、BPSK调制、四类信道支路、相应接收处理、硬判决、BCH译码和整帧统计之间的关系。

\begin{table}[htbp]
  \centering
  \scriptsize
  \caption{BCH多信道适应性仿真设置}
  \label{tab:bch-multichannel-settings}
  \setlength{\tabcolsep}{2pt}
  \begin{tabular}{p{0.14\textwidth}p{0.21\textwidth}p{0.26\textwidth}p{0.18\textwidth}p{0.08\textwidth}}
    \toprule
    信道场景 & 信道模型 & 关键参数 & 接收处理 & 主指标 \\
    \midrule
    固定多径 & 四抽头实数信道叠加AWGN & 归一化多径抽头：0.80452、0.52294、0、0.28158 & 已知信道，线性MMSE均衡 & FER \\
    帧内相位累积 & 复基带符号相位线性增长并叠加AWGN & 帧首$0^\circ$、帧末累计$30^\circ$ & 不作频偏补偿，取实部硬判决 & FER \\
    连续遮挡 & 连续调制符号幅度置零并叠加AWGN & 随机起点，约占码长10\% & 不提供受损位置，不交织 & FER \\
    连续突发错误 & 连续硬比特人为翻转 & 随机起点，0--50 bit，不环回 & 无AWGN，不交织 & FER \\
    \bottomrule
  \end{tabular}
\end{table}

\begin{figure}[htbp]
  \centering
  \fbox{\begin{minipage}[c][0.21\textheight][c]{0.82\textwidth}
    \centering\Large 待插入：Visio绘制的BCH多信道适应性仿真链路图
  \end{minipage}}
  \caption{BCH多信道适应性仿真链路}
  \label{fig:bch-multichannel-placeholder}
\end{figure}

\subsection{固定多径条件}

固定多径信道表示为$y_i=\sum_{\ell}h_\ell x_{i-\ell}+n_i$。四个多径抽头经过能量归一化，延迟副本之间的叠加造成码间干扰。接收端已知信道系数，采用线性最小均方误差（Minimum Mean Square Error，MMSE）均衡恢复符号，再作硬判决和BCH译码。

\reportfigure[0.96\textwidth]{figures/bch/s2_multipath_fer.png}
  {固定实数多径与MMSE均衡后的帧错误率}{fig:bch-s2-multipath}
  {横轴为统一换算后的$E_s/N_0$。}

图\ref{fig:bch-s2-multipath}中的曲线随信噪比提高而继续下降，说明已知信道条件下的MMSE均衡能够抑制主要码间干扰，没有形成明显不可消除的平台。与AWGN基准相比，下降区间整体后移，其原因是有限长度块均衡仍保留边界失真，同时MMSE权衡会带来残余干扰和噪声增强。各方案的相对趋势总体保持：较大的整块纠错能力更能吸收均衡后分散出现的硬判决错误。不过，该结论建立在固定实数多径抽头和完全已知信道上，不能直接外推到信道估计误差或快速时变环境。

\subsection{频偏引起的帧内相位累积（累计30°）}

频偏会使载波相位随符号时刻持续旋转。本试验不施加固定的初始相差，而是令每个编码帧从$0^\circ$开始，相位按符号序号线性增加，并在帧末累计到$30^\circ$；因此，“累计30°”描述的是一个编码帧内由频偏引起的相位增长强度。不同码长方案均在各自帧首和帧末之间完成相同的总相位累积，接收端关闭频偏补偿，只取复基带样值实部进行硬判决。

\reportfigure[0.96\textwidth]{figures/bch/s2_phase_drift_fer.png}
  {频偏引起的帧内相位累积（累计30°）条件下的帧错误率}{fig:bch-s2-phase}
  {帧首相位为$0^\circ$，帧末相位为$30^\circ$，接收端不作补偿。}

对理想BPSK符号$\pm1$，旋转后的实部分量为$\pm\cos\theta$。在帧末$\theta=30^\circ$时，$\cos30^\circ\approx0.866$，判决幅度和噪声裕量有所减小，但符号能量并未像遮挡那样被直接删除。因此，随着信噪比提高，各方案的FER仍能继续下降，没有出现与10\%连续置零相同的平台。较强整块码对相位累积造成的分散判决错误仍表现出更好的吸收能力，分组方案则更容易因某个短子块发生多错而产生整帧错误。这里评价的是编码在未补偿相位累积下的承受能力，并非载波同步算法本身的性能。

\subsection{连续遮挡条件}

遮挡试验从随机位置选择约占编码帧10\%的连续调制符号，将其幅度直接置零，遮挡区仍叠加噪声。接收端既不知道受损位置，也没有擦除标记或交织处理。因此，该模型不是擦除码信道，而是“局部信号缺失后仍按普通硬判决输入BCH译码器”的压力试验。

\reportfigure[0.96\textwidth]{figures/bch/s2_blockage_fer.png}
  {约10\%连续符号置零遮挡下的帧错误率}{fig:bch-s2-blockage}
  {遮挡起点随机，受损位置不提供给译码器。}

图\ref{fig:bch-s2-blockage}表现出显著错误平台。提高AWGN信噪比只能减小未遮挡区域的噪声，却无法重建被置零的连续符号；当局部损伤产生的错误数超过BCH码纠错能力时，继续增加信噪比也难以改善整帧结果。BCH(511,385)整块方案凭借$t=14$取得相对较低的FER，但仍不能消除平台。该现象与累计30°相位旋转形成清楚对比：相位旋转主要压缩实部判决裕量，连续遮挡则直接丢失一段符号信息。对于强遮挡业务，必须考虑交织、受损位置辅助或更有针对性的接收结构。

\subsection{连续突发错误条件}

连续突发错误试验不叠加AWGN，而是在硬判决比特序列中从随机起点连续翻转0--50 bit，不允许越过帧尾环回，也不采用交织。它是为了隔离“错误集中程度”而设置的人工模型，不能等同于自然无线信道中的独立随机误码。

\reportfigure[0.96\textwidth]{figures/bch/s2_burst_fer.png}
  {无AWGN、无交织条件下的连续突发错误帧错误率}{fig:bch-s2-burst}
  {横轴为人为施加的连续翻转长度。}

当单个码字内的翻转数不超过$t$时，BCH译码具有确定的纠错保证；超过$t$后，是否恢复成功还取决于突发起点、码字边界以及错误是否跨入相邻子块。分组方案的每个15 bit子块只能保证纠正1 bit，连续突发很容易在某一子块内超过能力。双块BCH(255,207)方案则应按每个子块分别判断：跨越两块边界的突发可能把错误分摊，但整帧仍要求两块都正确，不能把总突发长度简单等同于某一个$t$值。图中曲线因此是特定边界分布下的统计结果，而不是“突发长度等于纠错能力”的普遍规则。

\subsection{不同错误形态对BCH性能的影响}

从错误空间分布看，本章信道可归为两类。第一类以分散错误为主，包括AWGN硬判决错误、固定多径均衡后的残余错误以及累计30°相位旋转下的判决错误。这些错误可以分布在较长码字范围内，提高$t$通常能够使整块BCH方案吸收更多错误，FER曲线也会随信噪比持续下降。多径中的残余码间干扰可能带有相关性，相位累积又使帧后部的判决裕量更小，但二者在当前条件下仍未表现为连续大段信息完全丢失。

第二类是局部集中的连续损伤，包括调制符号置零遮挡和人为连续翻转。它们可能一次向同一子码字注入超过$t$的错误，因而仅靠增加信噪比或适度增强BCH码不足以保证恢复。改善方向包括在编码前交织以重新分散错误、向译码器提供受损位置、利用导频或幅度检测识别遮挡区，以及设计面向连续损伤的接收处理。相关结构将在第\ref{chap:interleaving}章结合交织进一步讨论。本章不建立跨不同横轴和不同损伤机理的统一排名，只在相同信道、相同判据下比较方案。

\section{BCH方案阶段性结论}

表\ref{tab:bch-conclusions}把仿真结果转化为可执行的工程建议。推荐项均以本章已有数据为依据，不额外设定未经验证的工作门限。

\begin{table}[htbp]
  \centering
  \scriptsize
  \caption{BCH方案的阶段性工程建议}
  \label{tab:bch-conclusions}
  \setlength{\tabcolsep}{3pt}
  \begin{tabular}{p{0.17\textwidth}p{0.22\textwidth}p{0.35\textwidth}p{0.17\textwidth}}
    \toprule
    业务条件 & 推荐方案 & 推荐理由 & 主要代价 \\
    \midrule
    200 bit，低复杂度优先 & 分组BCH(15,11) & 短码查表路径简单，软件译码时间最低 & 发送285 bit；AWGN可靠性较弱 \\
    200 bit，可靠性优先 & BCH(511,385)整块方案 & AWGN曲线最靠左，集中损伤下相对表现最好 & 发送326 bit；软件译码时间最高 \\
    300 bit，综合平衡 & BCH(511,421)整块方案 & 390 bit码长下兼顾$t=10$、FER与软件时间 & 复杂度高于分组方案 \\
    300 bit，可靠性优先 & BCH(511,385)整块方案 & AWGN下FER最低，具有$t=14$纠错能力 & 比390 bit方案多发送36 bit \\
    强连续损伤 & 不依赖单一BCH方案 & 遮挡平台和突发结果表明局部错误可一次超过$t$ & 需增加交织、位置辅助或结构化接收处理 \\
    \bottomrule
  \end{tabular}
\end{table}

总体而言，BCH(15,11)分组方案适合强调实现简单和处理开销的低速链路；较长的缩短BCH码通过联合纠正更大范围内的错误，能够显著改善AWGN和其他分散错误条件下的整帧可靠性。对于300 bit电文，BCH(511,421)整块方案是较均衡的基准；对于可靠性优先的200 bit或300 bit业务，BCH(511,385)整块方案更具优势。若链路存在强连续遮挡或长突发错误，则不应把提高BCH纠错能力作为唯一措施，而应把编码与交织、同步、受损位置检测和接收机处理联合设计。

需要强调的是，上述建议建立在系统编码、BPSK硬判决以及本章给定停止条件之上。不同方案的实际码率并不相同，曲线差异同时包含冗余长度、纠错能力、单块或多块组织方式等因素，因此不能把任意两条曲线间距直接解释为与实现无关的固定编码增益。软件译码时间同样只用于当前实现的相对比较；若后续改用并行逻辑、流水线Chien搜索或专用有限域运算单元，复杂度排序可能保持，但绝对处理时间需要重新测量。

从系统应用角度看，BCH码适合作为低速控制电文的基础纠错手段：在分散噪声占主导时，可根据允许的发送长度选择$t=10$或$t=14$的整块方案；在处理资源受限时，可退回结构清晰的分组方案。对连续损伤占主导的链路，后续设计应优先改变错误进入译码器之前的空间形态，再讨论是否继续增加BCH冗余。这样能够避免用过长码字换取有限改善，也使编码选择与实际信道机理保持一致。
